\documentclass{beamer}

\usetheme{Boadilla}
\usepackage{polski}
\usepackage{times}
\usepackage[T1]{fontenc}
\usepackage{graphicx}
\usepackage{listings}
\usepackage{hyperref}
\usepackage{subcaption}
\usefonttheme{serif}

\title{Seminarium dyplomowe, 3}
\subtitle{Układ aktywnego tłumienia hałasu pod nadzorem jądra systemu Linux z rozszerzeniami real-time}

\author{Daniel Kopiecki}
\institute{Politechnika Śląska}
\date{\today}

\begin{document}

\begin{frame}
\titlepage
\end{frame}

\begin{frame}
\frametitle{Słowem wstępu - system operacyjny}
\begin{itemize}
\item System operacyjny w najprostszej definicji to oprogramowanie zarządzające systemem komputerowym.
\item Z punktu widzenia użytkownika, stanowi środowisko, pozwalające na interakcję z komputerem
bez zrozumienia i znajomości działania urządzenia na poziomie sprzętowym.
\item W uogólnieniu system operacyjny zajmuje się:
\begin{enumerate}
    \item Planowanie i przydziałem czasu procesora na wykonanie poszczególnych zadań (programów).
    \item Kontrolą i przydziałem pamięci na zadania (programy).
    \item Synchronizacją zadań (programów).
    \item Obsługą zasobów sprzętowych, zapewnieniem jednolitego interfejsu kontroli nad nim.
\end{enumerate}
\end{itemize}
\end{frame}

\begin{frame}
\frametitle{Słowem wstępu - system operacyjny}
\begin{itemize}
    \item Przyjęto, że system operacyjny dzieli się na trzy elementy:
    \begin{itemize}
        \item \textbf{Jądro} - kontrola zadań i zasobów (sprzęt, pamięć, czas procesora).
        \item \textbf{Powłoka} - interakcja użytkownika z jądrem.
        \item \textbf{System plików} - sposób strukturyzacji danych na nośniku.
    \end{itemize}
\end{itemize}
\begin{figure}[H]
    \begin{minipage}{.3\textwidth}
      \centering
      \includegraphics[width=\linewidth]{../images/user.png}
    \end{minipage}%
    \begin{minipage}{.3\textwidth}
      \centering
      \includegraphics[width=\linewidth]{../images/os.png}
    \end{minipage}
    \caption{System operacyjny}
\end{figure}
\end{frame}

\begin{frame}
\frametitle{Jądro systemu operacyjnego}
\begin{itemize}
    \item Jądro systemu operacyjnego to program, który posiada pełną kontrolę
    nad wszystkim co dzieje się w systemie komputerowym.
    \item Kod tego programu jest umiejscowiony w odpowiednim, chronionym przed zapisem
    obszarze pamięci, a ładowany jest przy uruchomieniu systemu z użyciem \textbf{bootloader'a}.
    \item Istnieją różne architektury jądra, z których wyróżnić można:
    \begin{itemize}
        \item \textbf{Monolityczne} - jądro stanowi jeden, duży program (systemy uniksopodobne, lecz nie tylko) pracujący w \textbf{trybie jądra}.
        \item Mikrojądro - monolit zostaje podzielony na mniejsze moduły, z których część pracuje w \textbf{trybie użytkownika} komunikujac się z serwerami (pośrednikami).
        \item Nanojądro - jeszcze większe podzielenie na moduły.
        \item Hybrydowe - kompromis pomiędzy monolitem, a modularyzacją. 
    \end{itemize}

    
\end{itemize}
\end{frame}

\begin{frame}
\frametitle{Jądro systemu operacyjnego}
\begin{figure}[H]
    \centering
    \includegraphics[width=\linewidth]{../images/kernels.png}
    \caption{Schematy interakcji}
\end{figure}
\begin{itemize}
    \item Monolit - jądro UNIX (np. BSD), \textbf{Linux}, MS-DOS
    \item Mikro/nano - QNX, Mach
    \item Hybryda - Windows NT/2000/XP/Vista/7/8/10/11
\end{itemize}
\end{frame}

\begin{frame}
\frametitle{Jądro monolityczne}
\begin{itemize}
    \item Cechuje się wysoką wydajnością, ponieważ jedynym pośrednikiem między zasobami,
    a aplikacjami, jest właśnie jądro.
    \item Jądro zapewnia względnie wysokopoziomowe API do interakcji z zasobami (właściwie samoistnie je definiuje).
    \item Interakcje z jądrem odbywają się poprzez wywoływania systemowe (ang. \textit{system call}).
\begin{figure}[H]
    \centering
    \includegraphics[width=0.3\linewidth]{../images/kernel.png}
    \caption{Schemat interakcji - jądro monolityczne}
\end{figure}
\item Dodatkową możliwością w tego typu architekturze jest dodawanie sterowników jądra jako ładowalnych modułów (ang. \textit{loadable kernel modules}).
\end{itemize}
\end{frame}

\begin{frame}
\frametitle{Łamanie trybu użytkownika (\textit{breaking user's space})}
\begin{itemize}
    \item Rozwiązaniem krytykowanym przez dużą rzeszę programistów jest 
    oddzielanie kwestii systemowych do pracy w trybie użytkownika (jądra hybrydowe).
    \item Tego typu rozwiązanie nie dość, że powoduje przeogromne problemy z zachowaniem 
    kompatybilności wstecznej, to jeszcze całkowicie uzależnia programistów pracujących pod
    danym systemem operacyjnym od jego wersji.
    \item W słowach Linusa Torvaldsa (LKML)
    \\
    \par \textit{Breaking user programs simply isn't acceptable. (…) We know that people use old binaries for years and years, and that making a new release doesn't mean that you can just throw that out. You can trust us.}
    \item Przykładem może być stara gra komputerowa (zazwyczaj zaprojektowana do działania pod pewnym systemem operacyjnym), której uruchomienie dzisiaj jest wyjątkowo trudne z powodu braku kompatybilności.
\end{itemize}
\end{frame}

\begin{frame}
\frametitle{POSIX}
\begin{itemize}
\item POSIX (ang. \textit{Portable Operating System Interface for UNIX}, IEEE 1003) to standard oprogramowania systemu operacyjnego, w który skład wchodzi:
\begin{itemize}
    \item Interfejs programistyczny (\textit{<unistd.h>}, read, write, fork, pipe...).
    \item Interfejs użytkownika, przykładowo polecenia powłoki \textit{echo, cat, grep}.
    \item Właściwości powłoki.
\end{itemize}
\item Pojęcie powstało na wskutek prób standaryzacji różnych odmian systemu operacyjnego UNIX. 
\begin{figure}[H]
    \centering
    \includegraphics[width=0.5\linewidth]{../images/unix.jpg}
\end{figure}

\end{itemize}
\end{frame}

\begin{frame}
\frametitle{Jądro Linux}
\begin{itemize}
    \item Linux to jądro systemu operacyjnego stworzone początkowo przez Linusa Torvaldsa w 1991 roku.
    \item Jest oparte na licencji GNU General Public License v2 i początkowo stanowiło wolny odpowiednik
    systemu UNIX.
    \item Jest jądrem monolitycznym i uniksopodobnym (jest w większości zgodne ze standardem POSIX).
\begin{figure}[H]
    \begin{minipage}{.2\textwidth}
      \centering
      \includegraphics[width=\linewidth]{../images/tux.png}
    \end{minipage}%
    \begin{minipage}{.4\textwidth}
      \centering
      \includegraphics[width=\linewidth]{../images/linux.png}
    \end{minipage}
    \caption{Jądro systemu Linux}
\end{figure}
    \item W połączeniu z narzędziami i bibliotekami systemowymi z projektu GNU, otrzymuje się
    kompletny system operacyjny \textbf{GNU/Linux}.
\end{itemize}
\end{frame}

\begin{frame}
\frametitle{Wbudowany GNU/Linux}
\begin{itemize}
    \item Z powodu wszechstronności i otwartych kodów źródłowych, system operacyjny
    GNU/Linux jest szeroko wykorzystywany w systemach wbudowanych.
    \item Kod źródłowy może być portowany na różne architektury.
    \item Z racji wszechostronności systemu system
    można modyfikować usuwając i dodawając moduły.
    \item Pozwala to na tworzenie dedykowanych systemów operacyjnych stworzonych
    na potrzeby tworzonych urządzeń.
\end{itemize}
\begin{figure}[H]
    \centering
    \includegraphics[width=0.5\linewidth]{../images/embedded_tux.png}
\end{figure}
\end{frame}

\begin{frame}
\frametitle{Wbudowany GNU/Linux - ograniczenia}
\begin{itemize}
    \item System GNU/Linux nie jest systemem czasu rzeczywistego.
    \item Jest to znaczące ograniczenie, przeszkadzające w wykorzystaniu systemu
    do urządzeń realizujących zadania krytyczne, gdzie wymagany jest w pełni
    deterministyczny czas realizacji.
    \item Powstało jednak sporo projektów skupiających się na przybliżeniu/dostosowaniu systemu GNU/Linux
    do systemu czasu rzeczywistego.
    \item Jednym z najczęściej wykorzystywanych rozwiązań jest projekt \textbf{Real-Time Linux}, znany
    ściślej jako zestaw łatek (ang. \textit{patch}) jądra o nazwie \textbf{PREEMPT\_RT}.
\begin{figure}[H]
    \centering
    \includegraphics[width=0.2\linewidth]{../images/rt-linux.png}
\end{figure}

\end{itemize}
\end{frame}

\begin{frame}
\frametitle{System czasu rzeczywistego, RTC (ang. \textit{Real-time computing})}
\begin{itemize}
    \item Aby można stwierdzić było, że system jest systemem czasu rzeczywistego, poza logiczną prawidłowością 
    wykonanych akcji, musi on również gwarantować określony czas ich realizacji.
    \item Zważając na określony czas realzacji, systemy czasu rzeczywistego przyjęło się dzielić na:
    \begin{itemize}
        \item Hard real-time - przekroczenie limitu czasowego traktowane jest jak awaria systemu.
        \item Firm real-time - przekroczenie limitu czasowego skutkuje dyskredytacją wykonywanego zadania.
        \item Soft real-time - przekroczenie limitu czasowego skutkuje obniżeniem priorytetu wykonywanego zadania (ang. \textit{usefulness of a result degrades}).
    \end{itemize}
    \item W przypadku systemów jednozadaniowych sprawa jest trywialna, jednak komplikuje się, 
    kiedy w grę wchodzi wielozadaniowość.
    \item Systemy operacyjne są wielozadaniowe.
\end{itemize}
\end{frame}

\begin{frame}
\frametitle{Planowanie i wywłaszczenie}
\begin{itemize}
    \item Jądro systemu Linux implementuje \textbf{planistę} (ang. \textit{scheduler}), którego zadaniem jest
    szeregowanie zadań w systemie (procesy, wątki).
    \item Tryb pracy planisty w jądrze Linux może być wybrany przez użytkownika. Domyślnym jest CFS (ang. \textit{Completely Fair Scheduler})
    i jego wybranie wynika z deklaracji zadania jako SCHED\_NORMAL.
    \item Zadania są uruchamiane na podstawie ich priorytetu. Planista iteruje po obsługiwanych zadaniach i
    w pierwszej kolejności uruchamia te z największym priorytetem.
    \item Jądro systemu Linux zapewnia \textbf{wywłaszczenie zadania} z poziomu przestrzeni użytkownika.
\end{itemize}
\end{frame}

\begin{frame}
\frametitle{Wywłaszczenie}
\begin{itemize}
    \item \textbf{Wywłaszczenie} (ang. \textit{preemption}) to akcja tymczasowego zatrzymania aktualnie wykonywanego zadania, mająca na celu
    wykonanie innego, zazwyczaj o większym priorytecie. 
    \item Po wykonaniu zadania przerywającego, następuje powrót do wczesniej
    wykonanywanego. Akcję zatrzymania nazywa się \textbf{przerwaniem}.
\end{itemize}
\begin{figure}[H]
    \centering
    \includegraphics[width=0.75\linewidth]{../images/preempt.png}
\end{figure}
\end{frame}

\begin{frame}
\frametitle{Wywłaszczenie}
\begin{itemize}
    \item W jądrze systemu Linux domyślnie wywłaszczone mogą zostać jedynie zadania pochodzące z przestrzeni
    użytkownika. 
    \item Oznacza to, że jeżeli przerwanie nadejdzie w trakcie wykonywania przez zadanie wywołania
    systemowego, to takie przerwanie nie zostanie obsłużone do momentu zakończenia zawołania systemowego.
    \begin{figure}[H]
    \centering
    \includegraphics[width=\linewidth]{../images/preempt_1.png}
\end{figure}
\end{itemize}
\end{frame}

\begin{frame}
\frametitle{Wywłaszczenie}
\begin{itemize}
\item Dodatkowo, każde zadanie z przestrzeni użytkownika może zostać wywłaszczone przez
przerwanie sprzętowe. W trakcie obsługi takiego przerwania zablokowana jest obsługa innych.
\item Oznacza to, że mało ważne przerwania sprzętowe mogą wywłaszczyć ważne, krytyczne czasowo
zadania czasu rzeczywistego.
    \begin{figure}[H]
    \centering
    \includegraphics[width=\linewidth]{../images/hard_irq.png}
    \end{figure}
\end{itemize}
\end{frame}

\begin{frame}
\frametitle{Wywłaszczenie}
\begin{itemize}
\item Jakby tego było mało, istnieje jeszcze szereg problemów z synchronicznym dostępem do zasobów.
\item Może zaistnieć sytuacja, kiedy niekrytyczne zadanie blokuje (ang. \textit{mutex}) dostęp do zasobów zadaniu krytycznemu.
    \begin{figure}[H]
    \centering
    \includegraphics[width=0.8\linewidth]{../images/tasks.png}
    \end{figure}
\item Krytyczne zadanie A musi uzyskać dostęp do zasobów blokowanych przez zadanie C (niski priorytet).
\item Planista więc uruchamia zadanie C, aby odblokować dostęp.
\item W tym czasie zadanie C zostaje wywłaszczone przez zadanie B o wyższym priorytecie, ale niekrytyczne.
\item Taki problem nosi nazwę \textbf{inwersji priorytetów} (ang. \textit{priority inversion}).
\end{itemize}
\end{frame}



\begin{frame}
\frametitle{Projekt PREEMPT\_RT}
\begin{itemize}
    \item Zestaw łatek PREEMPT\_RT celuje w możliwości wywłaszczenia wszystkich zadań realizowanych w systemie,
    nawet tych wykonywanych z przestrzeni jądra.
    \item Jako, że jądro Linux umożliwia korzystania z innych planistów niż domyślne, istnieje możliwość zastąpienia domyślnych
    wytycznych planisty na wspieranie rozszerzeń czasu rzeczywistego.
    \item Jako, że projekt PREEMPT\_RT oferuje jedynie zestaw łatek, interfejs programistyczny POSIX/Linux pozostaje 
    niezmienny.
    \item Umożliwa to korzystanie z systemu w tradycyjny sposób jak i wykonywanie krytycznych zadań równolegle w opraciu o 
    planistę z rozszerzeniami czasu rzeczywistego.
\end{itemize}
\end{frame}

\begin{frame}
\frametitle{Projekt PREEMPT\_RT}
\begin{itemize}
    \item Kilka funkcji łatkek PREEMPT\_RT:
    \begin{itemize}
        \item Wywłaszczenie zadań z przestrzeni jądra.
        \item Liczniki wysokiej rozdzielczości (ang. \textit{High resolution timers}).
        \item Wątkową realizację przerwań sprzętowych.
        \item Implementację blokad w trybie czasu rzeczywistego (do synchronizacji zadań).
        \item Planistę opartego o algorytm Earliest-Deadline First.
        \item Rozwiązanie problemu inwersji priorytetów poprzez zastosowanie \textbf{dziedziczenia priorytetu}.
    \end{itemize}
    \item Jako, że czas wykonania zadania jest bardzo trudny do predykcji, celem jest minimalizacja latencji pomiędzy reakcją na wykonanie zadania.
    \begin{figure}[H]
        \centering
        \includegraphics[width=0.6\linewidth]{../images/latency.png}
    \end{figure}
\end{itemize}
\end{frame}

\begin{frame}
\frametitle{Projekt PREEMPT\_RT}
\begin{itemize}
\item Wątkowa realizacja przerwań sprzętowych.
\begin{figure}[H]
    \centering
    \includegraphics[width=0.5\linewidth]{../images/threaded_irq.png}
    \end{figure}
\item Dziedziczenie priorytetu.
\begin{figure}[H]
\centering
\includegraphics[width=0.6\linewidth]{../images/prio_inh.png}
\end{figure}
\end{itemize}
Krytyczne zadanie A musi uzyskać dostęp do zasobów blokowanych przez zadanie C. W związku z tym,
priorytet zadania C jest zwiększany do priorytetu zadania A.
\end{frame}

\begin{frame}
\frametitle{Embedded GNU/Linux, PREEMPT\_RT - Proces uruchamiania, kolejność}
\begin{figure}[H]
    \centering
    \includegraphics[width=\linewidth]{../images/boot.jpg}
\end{figure}
\end{frame}

\begin{frame}[fragile]
\frametitle{Embedded GNU/Linux - przygotowywanie nośnika}
\begin{lstlisting}
sudo parted --script /dev/kartaSD mklabel msdos \
    mkpart primary fat32 0% 100MiB \
    mkpart primary ext4 100MiB 1100MiB \
    mkpart primary ext4 1100MiB 100% \
    set 1 boot on

sudo mkfs.vfat /dev/kartaSD1 -n uboot
sudo mkfs.ext4 /dev/kartaSD2 -L kernel
sudo mkfs.ext4 /dev/kartaSD3 -L rootfs
\end{lstlisting}
\end{frame}

\begin{frame}[fragile]
\frametitle{Embedded GNU/Linux, PREEMPT\_RT - bootloader}
\begin{itemize}
\item Das U-Boot (ang. \textit{the Universal Boot Loader}) - główny, otwartoźródłowy bootloader używany
w zastosowaniach embedded GNU/Linux.
\begin{figure}[H]
    \begin{minipage}{.3\textwidth}
      \centering
      \includegraphics[width=0.6\linewidth]{../images/uboot.png}
    \end{minipage}%
    \begin{minipage}{.3\textwidth}
      \centering
      \includegraphics[width=\linewidth]{../images/uboot_view.png}
    \end{minipage}
\end{figure}
\item Kompilacja i kofiguracja bootloader'a:
\begin{lstlisting}[basicstyle=\tiny]
$ git clone --depth 1 https://github.com/u-boot/u-boot.git
$ cd u-boot
$ make ARCH=arm CROSS_COMPILE=arm-linux-gnueabihf- am335x-evm-defconfig -j5
$ make ARCH=arm CROSS_COMPILE=arm-linux-gnueabihf- -j5
$ cp MLO /media/user/uboot/
$ cp u-boot.img /media/user/uboot/
\end{lstlisting}
\end{itemize}
\end{frame}

\begin{frame}[fragile]
\frametitle{Embedded GNU/Linux, PREEMPT\_RT - jądro i system plików}
\begin{itemize}
\item Kompilacja i konfiguracja jądra:
\begin{lstlisting}[basicstyle=\tiny]
    $ git clone --depth 1 https://github.com/beagleboard/linux.git 5.10.162-ti-rt-r59
    $ cd linux
    $ make ARCH=arm CROSS_COMPILE=arm-linux-gnueabihf- bb.org_defconfig -j5
    $ make ARCH=arm CROSS_COMPILE=arm-linux-gnueabihf- -j5
    $ cp arch/arm/boot/zImage /media/user/kernel
    $ cp arch/arm/boot/dts/ am335x-boneblack.dtb /media/user/kernel
\end{lstlisting}
\item Rozpakowanie gotowego systemu plików:
\begin{lstlisting}[basicstyle=\tiny]
    $ mkdir -p rootfs
    $ cd rootfs
    $ wget https://www.rcn-ee.com/rootfs/eewiki/minfs/ubuntu-20.04.5.tar.xz
    $ tar -xf ubuntu-20.04.5-minimal-armhf-2022-12-20.tar.xz
    $ rm ubuntu-20.04.5-minimal-armhf-2022-12-20.tar.xz
    $ cp -r * /media/user/rootfs
\end{lstlisting}
\end{itemize}
\end{frame}

\begin{frame}
\frametitle{Materiały źródłowe}
    \begin{itemize}
        \item https://www.linux.org/
        \item https://bootlin.com/
        \item https://archive.kernel.org/oldwiki/rt.wiki.kernel.org/
        \item https://wiki.linuxfoundation.org/realtime/start
    \end{itemize}
\vfill
    \large{Dziękuję za uwagę}
\end{frame}


\end{document}
